\documentclass[dvipdfmx]{jsarticle}
\usepackage{graphicx}
\usepackage{color}
\usepackage{url}
\usepackage{listings}
\usepackage{amsmath}
\usepackage{amssymb}
\usepackage{tikz}
\usetikzlibrary{shapes.geometric, arrows.meta, positioning, fit, backgrounds}

\begin{document}

\pagestyle{empty}

\begin{center}
\vspace*{1cm}
\large
{\LARGE 卒業研究報告書}\\
\vspace*{0.8cm}
題目\\
\vspace*{1cm}
{\Huge \underline{マルチモーダル学習による薬剤鑑別}}\\
\vspace{3mm}

\vspace*{3cm}
指導教員\\
\vspace*{0.3cm}
\underline{\LARGE 森山　真光　准教授}\\
\vspace*{3cm}
報告者\\
\vspace*{0.3cm}

{23--1--211--0009}\\
\vspace*{0.3cm}
\underline{\Huge 呉竹　櫂人}\\
\vspace*{0.5cm}
近畿大学情報学部情報学科\\
\vspace*{2cm}
令和7年6月2日提出\\
\end{center}

\newpage
\normalsize

%\tableofcontents
\clearpage
\begin{center}
{\bf \Large 概要}
\end{center}

薬剤の取り違えや誤服用は医療安全上の重要な課題であり，医療現場や家庭では，薬剤名が記載された外箱や薬袋が失われた場合，または薬剤本体やPTPシートのみが残された場合に，錠剤やカプセルの外観，刻印，包装上の印字など限られた視覚情報を手がかりに薬剤を鑑別する必要がある．本研究では，薬剤画像とテキストプロンプトを入力とし，薬剤名を出力するVision-Language Model型の薬剤鑑別手法を検討する．具体的には，オープンソースのマルチモーダルLLMであるQwen3.5を用い，薬剤画像に対して「この錠剤は何という薬ですか」という質問を与え，正解薬剤名を出力するようにLoRAによるファインチューニングを行う．比較対象として，ファインチューニング前のQwen3.5および外部API型マルチモーダルAIであるGeminiによるゼロショット識別を用い，薬剤画像データに特化したファインチューニングが薬剤鑑別精度に与える効果を評価する．41種類の薬剤を対象とした評価では，ファインチューニング後のモデルが単一薬剤識別でTop-1精度98.2\%を達成し，ファインチューニング前（0.9\%）およびGemini（82.3\%）を大きく上回った．さらに，1枚の画像に複数の薬剤が写る場合についても，薬剤位置を検出して領域ごとに単剤識別を行う2段構成を提案し，検出器の専用学習により完全一致率96.0\%を達成した．スマートフォン搭載を見据えたモデルサイズ（0.8B/2B/4B）および4bit量子化による精度変化の検証も行った．また，薬剤データベースを検索して候補薬剤情報を提示するRAG型手法については予備実験で比較を行っており，全薬剤規模での比較を今後の発展的課題とする．

\newpage

\tableofcontents
\newpage
\setcounter{page}{1}
\pagestyle{plain}

%----------------------------------------------------------------------
\section{序論}\label{sec1}
%----------------------------------------------------------------------

\subsection{本研究の背景}

薬剤の取り違えや誤服用は，患者の健康被害につながる医療安全上の重要な問題である．
世界保健機関（WHO）は，2017年に第3回世界患者安全チャレンジとして「Medication Without Harm」を立ち上げ，投薬に関連する重篤かつ回避可能な健康被害を5年間で50\%削減することを世界共通の目標として掲げている\cite{who_medication}．
WHOは，こうした投薬に関連する有害事象に世界全体で年間およそ420億ドルの費用が費やされていると報告しており，投薬の安全性確保は一国にとどまらない世界的な課題と位置づけられている．
国内においても，日本医療機能評価機構の薬局ヒヤリ・ハット事例収集・分析事業には毎年多数の調剤に関する事例が報告されており，「薬剤取り違え」は上位を占める調剤過誤の一つである\cite{jcqhc_hiyari}．
特に，名称や外観が類似した薬剤による取り違えが多く報告されており，薬剤を正確に識別・鑑別することは医療安全上の重要な課題である．

医療現場や家庭では，薬剤名が記載された外箱や薬袋が失われた場合，または薬剤本体やPTPシートのみが残された場合に，錠剤やカプセルの外観，刻印，包装上の印字など限られた視覚情報を手がかりに薬剤を鑑別する必要がある．
特に，高齢者が複数の薬剤を長期的に服用する場合や，一包化調剤によって複数の薬剤がまとめられる場合には，薬剤の識別が困難になり，誤服用のリスクが高まる．

従来，薬剤鑑別には薬剤師などの専門知識が必要であり，薬剤データベースを手作業で検索する負担も大きい．
実際に末次ら\cite{suetsugu2026ai}は，持参薬の一包化鑑別が薬剤師にとって時間を要する負担の大きい業務であることを報告している．
また，白色円形の錠剤のように外観が類似する薬剤も多く，画像のみから薬剤を正確に識別することは容易ではない．
こうした鑑別業務の負担を軽減する手段として，近年ではAIによる画像解析を用いた薬剤識別の自動化が注目されている．
同研究\cite{suetsugu2026ai}は，AI画像解析錠剤識別システムを持参薬の一包化鑑別業務に導入し，鑑別に要する時間を短縮できることに加え，従来は薬剤師が担っていた鑑別業務を薬剤助手へタスク・シフト／タスク・シェアできる可能性を報告している．
このことは，薬剤鑑別をモデルによって自動化することが，医療安全の確保と医療従事者の負担軽減の両面で有用であることを示している．
このような背景から，薬剤画像を用いて自動的に薬剤を識別する研究が行われてきた．

Heoら\cite{heo2023pill}は，薬剤識別を認識ステップと検索ステップに分けた深層学習ベースの薬剤識別システムを提案している．
YOLOv5により錠剤表面の刻印文字を検出し，ResNet-32により形状，色，剤形を認識する．
さらに，RNNベースの文字レベル言語モデルを用いて検出された刻印文字の誤りを補正し，薬剤データベース内の情報との類似度に基づいて候補薬剤を順位付けする．
この手法は，薬剤画像に含まれる視覚情報と刻印文字情報を組み合わせて識別する点で有効であり，検索型の構成を採用することで学習に含まれない薬剤への対応も考慮している．
一方で，複数の専用モジュールを組み合わせる必要があり，システム構成が複雑である．
また，刻印文字を重要な識別手がかりとしているため，刻印が不明瞭な画像や刻印のない薬剤では識別性能が低下する可能性がある．

Rádliら\cite{radli2024metric}は，薬剤画像の視覚特徴と薬剤説明テキストから得られるテキスト特徴を組み合わせたマルチモーダルな薬剤認識手法を提案している．
この研究は，画像情報だけでなくテキスト情報を利用して薬剤認識を行う点で本研究と近い．
一方で，同手法は画像特徴抽出器，テキスト埋め込み，メトリクス学習を個別に設計する構成であり，画像とテキストプロンプトを入力として薬剤名を直接生成するVision-Language Model型の手法ではない．

Menahaら\cite{gemini2024pill}は，Google Geminiを用いた生成AIベースの薬剤識別支援システムを提案している．
この研究は，マルチモーダルAIを薬剤識別支援に応用している点で本研究と関連する．
一方で，同研究は商用クローズドソースモデルを用いたプロンプトベースの手法であり，薬剤画像データを用いてモデル自体をファインチューニングする構成ではない．
また，外部APIを利用する場合，薬剤画像を外部サービスに送信する必要があり，プライバシー上の懸念が残る．

以上の課題から，（1）錠剤本体の外観，刻印，包装上の印字など複数の視覚情報を統合した識別，（2）外部APIに過度に依存しないローカル動作の可能性，（3）薬剤画像に特化したファインチューニングによる識別精度の向上，の3点を検討する必要がある．

\subsection{本研究の目的}

本研究では，オープンソースのVision-Language ModelであるQwen3.5を用い，薬剤画像とテキストプロンプトから薬剤名を出力する薬剤鑑別手法を検討する．
具体的には，薬剤画像に対して「この錠剤は何という薬ですか」という質問を与え，正解薬剤名を出力するようにLoRAによるファインチューニングを行う．

本研究の目的は，ファインチューニング前のQwen3.5および外部API型マルチモーダルAIであるGeminiと比較することで，薬剤画像データに特化したファインチューニングが薬剤鑑別精度に与える効果を検証することである．
また，オープンソースモデルを用いることで，薬剤画像を外部APIに送信せず，適切な計算環境を用意すればローカル環境で推論できる薬剤鑑別システムの可能性を示すことを目指す．

なお，本研究におけるマルチモーダル学習とは，薬剤画像を視覚情報として入力し，自然言語による質問文をテキスト情報として与え，薬剤名を応答として生成するVision-Language Model型の学習を指す．
本報告時点では，テキスト入力は質問文を中心とするが，今後は識別コード，薬剤候補名，形状情報，効能情報などを追加し，より情報量の多い画像・テキスト統合型の鑑別へ拡張することを検討する．

\subsection{本報告書の構成}

\ref{sec2}章では関連研究および関連技術を整理し，本研究の位置づけを明確にする．
\ref{sec3}章では提案手法の詳細および実験計画を述べる．
\ref{sec4}章では実験結果と考察を示す．
\ref{sec5}章では結論と今後の課題を述べる．

\newpage

%----------------------------------------------------------------------
\section{関連研究・関連技術}\label{sec2}
%----------------------------------------------------------------------

本研究では以下の3本の先行研究を参考文献とし，各手法の概要と限界を整理する．

\subsection{画像特徴・刻印文字を用いた薬剤識別手法}

Heoら\cite{heo2023pill}は，薬剤識別を認識ステップと検索ステップに分けた深層学習ベースの薬剤識別システムを提案した．
認識ステップでは，YOLOv5により錠剤表面の刻印文字を検出し，ResNet-32により形状，色，剤形を認識する．
さらに，RNNベースの文字レベル言語モデルを用いて，検出された刻印文字の誤りを補正する．
検索ステップでは，認識された形状，色，剤形，刻印文字と薬剤データベース内の情報との類似度を計算し，候補薬剤を順位付けする．
韓国の薬剤データセット（MFDS）においてTop-1精度85.6\%，米国データセット（NLM）において74.5\%を達成している．

この手法は，薬剤画像に含まれる視覚情報と刻印文字情報を組み合わせて識別する点で有効であり，画像分類ではなく検索型の構成を採用することで，学習に含まれない薬剤への対応も考慮している．

\textbf{課題}
\begin{itemize}
  \item YOLOv5，ResNet，RNNベースの言語モデルなど複数のモジュールを組み合わせる必要があり，システム構成が複雑である．
  \item 刻印文字を重要な識別手がかりとしているため，刻印が不明瞭な画像や刻印のない薬剤では，識別性能が低下する可能性がある．
  \item 検索型の構成により未知薬剤への対応を考慮している一方で，刻印検出や特徴認識の各モジュールが対象データに依存する可能性がある．
\end{itemize}

\subsection{画像特徴とテキスト特徴を用いたマルチモーダル薬剤識別手法}

Rádliら\cite{radli2024metric}は，薬剤画像の視覚特徴と薬剤説明文から得られるテキスト特徴を組み合わせたマルチモーダルな薬剤認識手法を提案している．
同研究では，EfficientNetV2をバックボーンとし，RGB画像，局所特徴，輪郭，テクスチャなど複数の視覚特徴を抽出する．
また，薬剤説明文から得られる単語埋め込みを用いて，メトリクス学習における動的マージントリプレット損失（DMTL）のマージンを決定する．
112クラス・4,480枚で構成するOGYEIv2データセットにてTop-1精度93.56\%を達成している．

この研究は，画像情報だけでなくテキスト情報を利用して薬剤認識を行う点で本研究と近い．

\textbf{課題}
\begin{itemize}
  \item 画像特徴抽出器，テキスト埋め込み，メトリクス学習を個別に設計する必要がある．
  \item 画像とテキストプロンプトから薬剤名を直接出力するVLM型の手法ではない．
  \item 評価データが特定地域・特定データセットに限定されており，他の薬剤画像への適用可能性は追加検証が必要である．
\end{itemize}

\subsection{生成AI・LLMを用いた薬剤識別支援手法}

Menahaら\cite{gemini2024pill}は，Google Geminiを用いた生成AIベースの薬剤識別支援システムを提案している．
同研究は，薬剤画像を入力として，薬剤に関する情報をテキストおよび音声で提示するシステムであり，マルチモーダルAIを薬剤識別支援に応用している点で本研究と関連する．

一方で，同研究は商用クローズドソースモデルを用いたプロンプトベースの手法である．
少なくとも同研究では，薬剤画像データを用いてモデル自体をファインチューニングする構成は示されていない．
また，外部APIを利用する場合，薬剤画像を外部サービスに送信する必要があり，プライバシー上の懸念が残る．
本研究では，オープンソースモデルであるQwen3.5を用い，ローカル環境でLoRAファインチューニングを行う点に差異がある．

\textbf{課題}
\begin{itemize}
  \item 商用クローズドソースモデルであり，当該研究ではモデル自体のファインチューニングは行われていない．
  \item 外部APIを利用する場合，薬剤画像を外部サービスに送信する必要があり，プライバシー上の懸念が残る．
  \item 利用コストが高く，継続的な運用において制約が生じる可能性がある．
  \item プロンプトベースの手法であり，当該研究では対象薬剤への追加学習は示されていない．
\end{itemize}

\subsection{先行研究の整理と本研究の位置づけ}

先行研究では，薬剤画像のみを用いるCNN分類に加え，刻印文字，薬剤説明テキスト，処方情報などを組み合わせるマルチモーダルな薬剤識別手法が提案されている．
Heoらの手法は，画像特徴と刻印文字を個別に抽出し，言語モデルで刻印文字を補正することで高精度な薬剤検索を実現している．
一方で，複数の専用モジュールを組み合わせる必要があり，システム構成が複雑である．
Rádliらの手法は，薬剤画像と薬剤説明テキストを組み合わせる点で本研究と関連するが，画像特徴抽出器とテキスト埋め込みを用いたメトリクス学習であり，薬剤名を直接生成するVLM型の手法ではない．
Menahaらの手法は，生成AIを薬剤識別支援に用いる点で関連するが，商用モデルを用いたプロンプトベースの手法であり，薬剤画像に特化したファインチューニングを行うものではない．

これらを踏まえ，本研究ではオープンソースのVision-Language ModelであるQwen3.5を用い，薬剤画像とテキストプロンプトを入力として薬剤名を直接出力する手法を検討する．
特に，LoRAによるファインチューニングを行うことで，薬剤画像データに特化したモデルを構築し，ファインチューニング前のモデルおよび外部API型マルチモーダルAIとの比較を通じて，その有効性を評価する．

表\ref{tab:related}に，先行研究の問題点と本研究での対応を整理して示す．

\begin{table}[h]
\centering
\caption{従来手法の問題点と本研究の対応}
\label{tab:related}
\small
\setlength{\tabcolsep}{4pt}
\begin{tabular}{p{0.22\linewidth}p{0.18\linewidth}p{0.52\linewidth}}
\hline
従来手法の問題点 & 該当研究 & 本研究での対応 \\
\hline
YOLOv5・ResNet・RNNなど複数の専用モジュールが必要でシステム構成が複雑 &
Heoら \cite{heo2023pill} &
Qwen3.5単一モデルで画像とテキストプロンプトを統合処理し，構成をシンプルに保つ \\
\hline
画像特徴抽出器・テキスト埋め込み・メトリクス学習を個別に設計する必要があり，エンドツーエンドの学習ではない &
Rádliら \cite{radli2024metric} &
LoRAによりQwen3.5全体を薬剤画像データに適応させ，エンドツーエンドでファインチューニングする \\
\hline
商用クローズドソースモデルへの依存および薬剤画像の外部送信によるプライバシー懸念 &
Menahaら \cite{gemini2024pill} &
オープンソースモデルをローカル環境で動作させることで，画像を外部に送信しない構成を目指す \\
\hline
ゼロショットのみで薬剤画像に特化した追加学習が行われておらず，専門的な薬剤識別への適応が不十分 &
Menahaら \cite{gemini2024pill} &
薬剤画像と正解薬剤名のペアを用いてLoRAファインチューニングを行い，識別精度を向上させる \\
\hline
\end{tabular}
\end{table}

\newpage

%----------------------------------------------------------------------
\section{提案手法}\label{sec3}
%----------------------------------------------------------------------

本章では，提案手法の概要と，課題・解決策・評価方法の対応関係を整理したうえで，手法の詳細を述べる．

\subsection{研究テーマの整理}

本研究の課題，解決策，評価方法の対応を表\ref{tab:theme}に示す．

\begin{table}[h]
\centering
\caption{研究テーマの課題・解決策・評価方法の対応}
\label{tab:theme}
\small
\setlength{\tabcolsep}{4pt}
\begin{tabular}{p{0.28\linewidth}p{0.38\linewidth}p{0.26\linewidth}}
\hline
課題 & 解決策 & 評価方法 \\
\hline
薬剤名が記載された外箱・薬袋が失われた場合に，錠剤・PTPシートの外観のみから薬剤を識別する必要がある &
オープンソースVLM（Qwen3.5-4B）を少数の実画像にデータ拡張を施した小規模データセットでLoRAファインチューニングし，単一モデルで薬剤鑑別を実現する &
Top-1 Accuracy（学習・評価データ分離）でGeminiのゼロショット識別と比較する \\
\hline
従来手法は複数の専用モジュールを要するか，商用APIに依存しプライバシー懸念がある &
オープンソースモデルをローカル環境で動作させることで，薬剤画像を外部サービスに送信しない構成を目指す &
外部API不使用でのローカル推論が可能かを動作確認する \\
\hline
\end{tabular}
\end{table}

\subsection{解決策の妥当性}

提案手法の妥当性を以下の観点から検討する．

\begin{itemize}
  \item \textbf{少数画像でのLoRA有効性}：単一薬剤（メインテート錠2.5mg）を元画像6枚から11倍に拡張した66枚で学習し，評価用22枚すべてで正答した（Top-1精度100\%）．少数の実画像でもデータ拡張との組み合わせによりLoRAが有効に機能することを予備実験で確認している．
  \item \textbf{多クラス拡張性}：41種類の薬剤を対象とした単一モデルのLoRAファインチューニングにおいても，評価用814枚中799枚で正答（全体Top-1精度98.2\%）を達成しており，薬剤数を増やした場合でも精度を維持できることを確認している．
  \item \textbf{外部APIとの比較根拠}：Geminiによるゼロショット識別との比較において，LoRAファインチューニング後のQwen3.5が単一薬剤識別（98.2\%対82.3\%）・複数薬剤同時識別（68.8\%対50.6\%）の双方で上回る結果を得ており，特化学習の有効性が示されている（\ref{sec4}章参照）．
  \item \textbf{複数薬剤への拡張可能性}：単剤識別用に学習したモデルをそのまま用い，薬剤位置の検出と領域ごとの単剤識別を組み合わせる2段構成（\ref{sec:multi}節）により，再学習なしで複数薬剤の同時識別が可能であることを確認している．さらに検出器を専用学習することで完全一致率96.0\%に到達した．
  \item \textbf{残課題}：回転（±30度）・水平反転に対する頑健性に課題があり，学習側の拡張強化または向き揃えの前処理を今後検討する．
\end{itemize}

\subsection{使用モデル：Qwen3.5}

本研究では，画像とテキストを入力できるオープンソースのVision-Language ModelとしてQwen3.5を使用する．
Qwen3.5を選定する理由は，外部APIに依存せずローカル環境で動作させられる可能性があり，LoRAなどのパラメータ効率の高い手法によって追加学習が可能であるためである．

Vision-Language Modelは画像とテキストプロンプトを同時に入力できるため，薬剤画像に対する質問応答形式で薬剤名を出力するタスクを構成できる．
本研究では，この特性を利用し，薬剤画像から薬剤名を出力するようにQwen3.5をLoRAによりファインチューニングする．

\subsection{提案手法：Qwen3.5 の LoRA ファインチューニング}

本研究では，薬剤画像とテキストプロンプトを入力し，正解薬剤名を出力するようにQwen3.5をLoRAによりファインチューニングする．
学習データは，薬剤ごとに分類された画像フォルダから作成する．
各サブフォルダ名を薬剤名ラベルとし，フォルダ内の画像を入力画像として用いる．

入力プロンプトは「この錠剤は何という薬ですか？薬剤名だけで答えてください。」とし，モデルが正解薬剤名のみを出力するように学習させる．
これにより，薬剤画像に対する質問応答形式の薬剤鑑別モデルを構築する．

\subsubsection{LoRAによる追加学習}

LoRA（Low-Rank Adaptation）は，事前学習済みモデルの重みをそのまま固定し，小さな追加パラメータ（低ランク行列）だけを学習するパラメータ効率的なファインチューニング手法である．
モデル全体を学習し直す場合に比べて学習対象のパラメータ数が大幅に少なくて済むため，少ない計算資源でも追加学習が可能になる．
本研究では，追加学習の大きさを決めるランクを$r=8$，スケーリング係数を$\alpha=16$とし，モデル内部のアテンション機構の問合せ射影（\texttt{q\_proj}）と値射影（\texttt{v\_proj}）にLoRAを適用する．

学習では，薬剤画像と質問プロンプトを入力として与え，正解薬剤名を出力するようにモデルを訓練する．
このとき，入力した質問プロンプトの部分は学習の対象とせず，モデルが答えるべき正解薬剤名の部分だけを学習させる．
これにより，モデルは画像中の薬剤に対して薬剤名のみを答えるように学習される．

図\ref{fig:pipeline}に，提案手法の処理フローを示す．

\begin{figure}[h]
\centering
\begin{tikzpicture}[
  box/.style={draw, rounded corners=3pt, minimum width=2.7cm, minimum height=1.1cm, align=center, font=\small, fill=white},
  model/.style={draw, rounded corners=3pt, minimum width=3.2cm, minimum height=1.8cm, align=center, font=\small, fill=blue!8},
  arrow/.style={-Stealth, thick}
]

% 左側：入力（薬剤画像と質問文）
\node[box] (img) {薬剤画像\\（PTPシート）};
\node[box, below=1.0cm of img] (prompt) {質問文\\「この錠剤は何\\という薬ですか？」};

% 中央：追加学習済みモデル（画像と質問文の中間の高さに配置）
\node[model, right=1.3cm of img, yshift=-1.05cm] (model) {Qwen3.5-4B\\＋\\LoRAアダプタ};

% 右側：出力（薬剤名）
\node[box, right=1.3cm of model] (out) {薬剤名を出力\\「メインテート錠」};

% 矢印
\draw[arrow] (img.east)    -- ++(0.5,0) |- (model.west);
\draw[arrow] (prompt.east) -- ++(0.5,0) |- (model.west);
\draw[arrow] (model.east)  -- (out.west);

\end{tikzpicture}
\caption{提案手法の処理フロー：薬剤画像と質問文をQwen3.5-4Bに入力し，LoRAアダプタで薬剤画像に適応させたモデルが薬剤名を出力する}
\label{fig:pipeline}
\end{figure}

\textbf{強み}
\begin{itemize}
  \item 薬剤画像データに特化した追加学習により，ゼロショット識別よりも高い識別精度が期待できる．
  \item 外部APIに依存せず，適切な計算環境を用意すればローカル環境で推論できる．
  \item 画像と自然言語プロンプトを組み合わせた質問応答形式で薬剤名を出力できる．
\end{itemize}

\textbf{課題}
\begin{itemize}
  \item ある程度の量と多様性を持つ学習データが必要である．
  \item 新たな薬剤クラスを追加する場合，再チューニングまたはデータベース検索との併用が必要となる．
  \item 包装込み画像を用いる場合，モデルが錠剤本体ではなく包装上の文字やデザインに依存する可能性がある．
\end{itemize}

\subsection{複数薬剤同時識別への拡張：検出→単剤分類}\label{sec:multi}

実際の薬剤鑑別では，1枚の画像に複数の薬剤（複数のPTPシート）が写る場合がある．
しかし，前節のモデルは「1枚に1薬剤」を答えるよう学習されているため，画像全体に対して「写っている薬をすべて答えよ」と質問しても1薬剤しか出力できない（この現象は\ref{sec4}章の実験で定量的に示す）．

そこで本研究では，再学習を行わずに複数薬剤を識別する2段構成のパイプラインを提案する．
基本的な考え方は，「複数薬剤の識別」を「薬剤位置の検出」と「領域ごとの単剤識別」に分解することである．
入力画像に対する処理手順を以下に示す．

\begin{enumerate}
  \item \textbf{検出}：Vision-Language Modelのグラウンディング能力を用いて，画像中の各薬剤シートの位置を四角（bounding box）で囲み，その座標をJSON形式で出力させる．検出プロンプトには「画像に写っているすべての薬剤(錠剤シート/PTPシート)の位置を検出し，各薬剤の bounding box を JSON で出力してください」を用いる．
  \item \textbf{切り出し}：検出した各四角を，画像の幅・高さの2\%のマージンを付けて切り出し，1薬剤ずつの部分画像にする．
  \item \textbf{単剤識別}：切り出した各部分画像に対して，\ref{sec3}.4節の単剤識別モデル（学習済みLoRA）を適用し，薬剤名を1つずつ得る．
  \item \textbf{統合}：すべての部分画像の識別結果をまとめ，その画像に写る薬剤名の集合として出力する．
\end{enumerate}

検出段階には2つの構成を用いる．
第1の構成では，LoRAを無効化した素のベースモデルのゼロショット検出能力を用いる（再学習・追加データ不要）．
第2の構成では，検出専用のLoRAアダプタを追加学習し，検出精度を高める（\ref{sec4}章で比較する）．
後者では，1つの基盤モデルに識別用LoRAと検出用LoRAの2つを搭載し，処理段階に応じてアダプタを切り替える．
このため，モデル本体のメモリ使用量はほぼ単剤識別時と変わらない．

図\ref{fig:multi_pipeline}に複数薬剤同時識別の処理フローを示す．

\begin{figure}[h]
\centering
\begin{tikzpicture}[
  box/.style={draw, rounded corners=3pt, minimum width=2.5cm, minimum height=1.5cm, align=center, font=\small, fill=white},
  stepbox/.style={draw, rounded corners=3pt, minimum width=2.5cm, minimum height=1.5cm, align=center, font=\small, fill=yellow!18},
  arrow/.style={-Stealth, thick}
]

% 上段：入力 →（1）検出 →（2）切り出し →（3）単剤識別
\node[box] (img) {複数薬剤画像\\（薬剤が\\複数枚）};
\node[stepbox, right=0.7cm of img] (det) {(1) 検出\\薬剤の位置を\\四角で出力};
\node[box, right=0.7cm of det] (crop) {(2) 切り出し\\各薬剤を\\1枚ずつに};
\node[stepbox, right=0.7cm of crop] (cls) {(3) 単剤識別\\各領域を\\個別に識別};

% 下段：（4）統合 → 出力
\node[box, below=0.9cm of cls] (merge) {(4) 統合\\答えを\\まとめる};
\node[box, left=0.7cm of merge] (out) {薬剤名の集合\\「メインテート，\\ウルソ」};

% 矢印（上段を右へ、折り返して下段を左へ）
\draw[arrow] (img)  -- (det);
\draw[arrow] (det)  -- (crop);
\draw[arrow] (crop) -- (cls);
\draw[arrow] (cls)  -- (merge);
\draw[arrow] (merge) -- (out);

\end{tikzpicture}
\caption{複数薬剤同時識別の処理フロー：検出→切り出し→単剤識別→統合の2段構成．検出（1）と単剤識別（3）は，同一のQwen3.5-4Bに検出用・識別用のLoRAアダプタを切り替えて実行する}
\label{fig:multi_pipeline}
\end{figure}

この方式の利点は，（1）高精度な単剤識別モデルを変更せずそのまま流用できること，（2）写っている薬剤の枚数が未知でも検出された個数から自動的に決まること，（3）検出段階のみを独立に改善できることである．

\subsection{比較手法}

提案手法の有効性を検証するため，以下の比較手法を設定する．

\begin{itemize}
  \item \textbf{ファインチューニング前のQwen3.5}：LoRA適用前のモデルによるゼロショット識別．
  \item \textbf{Gemini（\texttt{gemini-3.5-flash}）}：外部API型マルチモーダルAIによるゼロショット識別．
\end{itemize}

また，提案手法の出力形式を確定するための予備実験（\ref{sec4}章）では，上記に加えて以下の方式も比較した．

\begin{itemize}
  \item \textbf{詳細鑑別形式（方法B）}：薬剤名だけでなく，色・形状・刻印などの鑑別理由を含む詳細な説明を出力するようにファインチューニングする方式．
  \item \textbf{RAG型（薬剤データベース検索）}：薬剤情報を記述したデータベース（JSON形式）を検索し，候補薬剤情報を回答に利用する方式．
\end{itemize}

\subsection{使用データ}

本研究では，研究用に撮影した薬剤画像を使用する．
画像には，錠剤やカプセル本体に加え，PTPシートや包装上の印字が含まれる場合がある．
これは，実際の薬剤鑑別において，錠剤本体の色や形状だけでなく，包装上の文字情報や印字も重要な手がかりとなるためである．

ただし，包装情報を含む画像を用いる場合，モデルが錠剤本体ではなく包装上の文字やデザインに依存して薬剤名を予測する可能性がある．
本研究では，PTPシートや包装上の印字を含む画像をそのまま使用する．錠剤部分のみを切り出した画像との比較は今後の課題とする．

\subsection{評価方法}

評価においては，学習画像と評価画像が重複しないようにデータを分割する．
具体的には，各薬剤クラスごとに拡張前の元画像単位で学習用と評価用に分け，同一画像および同一の元画像に由来する拡張画像が学習データと評価データの両方に含まれないようにする．
これにより，単なる画像の記憶ではなく，薬剤画像に対する識別性能を評価する．
評価では，入力画像に対してモデルが出力した薬剤名と正解薬剤名を比較し，識別精度を算出する．

\subsubsection{単一薬剤識別の評価指標}

主な評価指標として，正答率（Top-1 Accuracy）を用いる．
正答率は，評価画像のうちモデルが正しく薬剤名を答えられた画像の割合であり，次式で定義される．

\begin{equation}
\text{正答率} = \frac{\text{正しく識別できた画像数}}{\text{評価画像の総数}}
\label{eq:top1}
\end{equation}

正解の判定は次のように行う．薬剤名は全角・半角や大文字・小文字といった表記のゆれを揃えたうえで，モデルの出力テキストの中に正解薬剤名が含まれていれば正解とみなす．
なお，後発医薬品のメーカー識別子（薬剤名末尾の「ＤＳＥＰ」のような「」内の表記）は，製品の鑑別に必須でなく画像からの判読も不確実であるため，正解薬剤名から除外して照合する（41薬剤中で該当するのはロスバスタチン錠2.5mgのみ）．この扱いはファインチューニング前後のQwen3.5・Geminiのすべてに統一して適用した．
また，モデルに上位3候補を出力させる実験では，3候補のいずれかに正解が含まれていれば正解とするTop-3正答率を補助指標として用いる．

\subsubsection{複数薬剤同時識別の評価指標}

複数薬剤の同時識別では，1枚の画像に写るすべての薬剤を当てる必要がある．
そこで主な指標として，予測した薬剤が正解の薬剤と過不足なく一致した画像の割合である完全一致率を用いる．

\begin{equation}
\text{完全一致率} = \frac{\text{予測が正解と完全に一致した画像数}}{\text{評価画像の総数}}
\label{eq:exact}
\end{equation}

完全一致率は，1薬剤でも取りこぼしや余分な予測があれば不正解とする厳しい指標である．
そこで補助指標として，部分的な正解にも部分点を与えるF1スコアも用いる．
F1スコアは，「予測した薬剤のうち正しかった割合（適合率）」と「正解の薬剤のうち予測できた割合（再現率）」をバランスよくまとめた指標であり，画像ごとに算出して全評価画像で平均する．

さらに，誤識別例については，外観が類似している薬剤，包装印字が不明瞭な薬剤，撮影条件が悪い画像などに分類し，誤識別の原因を分析する．

\subsection{実験計画}

本研究では，以下の手順で実験を進める．

\begin{enumerate}
  \item 予備実験として，単一薬剤を対象に出力形式（薬剤名のみ／詳細鑑別）およびRAG型・Geminiとの比較を行い，提案手法の出力形式を確定する．
  \item 薬剤画像データセットを作成し，薬剤名ラベルを整理する．
  \item 各薬剤クラスごとに，画像を学習用と評価用に分割する．この際，同一画像が学習データと評価データの両方に含まれないようにする．
  \item ファインチューニング前のQwen3.5を用いて，薬剤画像に対するゼロショット識別を行う．
  \item 薬剤画像と正解薬剤名のペアを用いて，Qwen3.5をLoRAによりファインチューニングする．
  \item ファインチューニング後のモデルを用いて同一条件で識別実験を行う．
  \item Gemini（\texttt{gemini-3.5-flash}）と比較する．
  \item 複数薬剤の同時識別を「検出→単剤分類」の2段構成（\ref{sec:multi}節）で評価し，ファインチューニング前のQwen3.5およびGeminiと同一の評価セットで比較する．
  \item 検出器の専用学習（検出用LoRA）により複数薬剤同時識別の改善を図り，検出が完璧な場合の上限値と比較する．
  \item スマートフォン搭載を見据え，4bit量子化およびモデルサイズ（0.8B / 2B / 4B）による精度の変化を，単一・複数・混在の各評価セットで測定する．
  \item 誤識別例を全実験横断で集計・分析し，薬剤外観，包装印字，撮影条件などが識別精度に与える影響を考察する．
\end{enumerate}

表\ref{tab:schedule}に実験スケジュールを示す．

\begin{table}[h]
\centering
\caption{実験スケジュール}
\label{tab:schedule}
\begin{tabular}{llp{0.55\linewidth}}
\hline
フェーズ & 期間 & 内容 \\
\hline
先行研究・準備 & 4〜5月 & 先行研究調査（論文3本以上），研究課題・解決策・評価方法の整理，薬剤画像撮影・データセット整備 \\
予備実験       & 4〜5月 & 単一薬剤LoRA予備学習（手法比較），全41薬剤拡張データセット学習 \\
本実験         & 5〜6月 & 単一薬剤3者比較，複数薬剤同時識別（検出→単剤分類・検出器学習），量子化・モデルサイズ比較，誤識別分析 \\
追加実験       & 7〜9月 & ネット収集画像での汎化評価，実写の複数薬剤検証，軽量2Bモデルの学習・評価 \\
論文執筆・発表 & 10月〜 & 卒業論文執筆，発表準備・スライド作成 \\
\hline
\end{tabular}
\end{table}

\newpage

%----------------------------------------------------------------------
\section{実験と結果}\label{sec4}
%----------------------------------------------------------------------

本章では，提案手法の有効性を検証するために実施した実験の内容と結果を述べる．
本研究は予備実験（単一薬剤での方式検証），本実験（41薬剤での単一・複数識別の3者比較），および追加実験（検出器の専用学習，軽量化評価）の順に進めた．
共同研究者による追実験を可能とするため，入力データの作成手順，処理方法，出力結果を，使用したパラメータとともに記述する．
本章に示すすべての数値は，実験実行時に自動出力された結果ファイル（評価サマリー・学習ログ）に基づく．

\subsection{実験環境}

実験に使用した計算環境を表\ref{tab:env}に示す．
学習・評価スクリプトはPythonで実装し，Hugging Face Transformers および PEFT ライブラリを用いた．
モデルの演算精度はbfloat16である．

\begin{table}[htbp]
\centering
\caption{実験環境}
\label{tab:env}
\begin{tabular}{ll}
\hline
項目 & 内容 \\
\hline
OS & Windows 11 \\
GPU & NVIDIA GeForce RTX 4070 Ti SUPER（VRAM 16GB） \\
Python & 3.10以上（uvによる仮想環境） \\
PyTorch & 2.10.0（CUDA 13.0） \\
transformers & 5.5.0 \\
peft & 0.18.1 \\
qwen-vl-utils & 0.0.14 \\
外部API & Gemini API（\texttt{gemini-3.5-flash}） \\
\hline
\end{tabular}
\end{table}

\subsection{データセットの構築}\label{sec:dataset}

\subsubsection{元画像の撮影}

調剤薬局で実際に取り扱われる薬剤を対象に，2026年4月に薬剤画像を撮影した．
各薬剤について，PTPシートの表面と裏面を含む3〜10枚程度の画像をスマートフォンのカメラで撮影し，薬剤名でフォルダを分けて保存した．
フォルダ名（正式な薬剤名・用量を含む）をそのまま正解ラベルとして用いる．
混合薬である一包化薬は単剤識別の対象外として除外した．
最終的なデータセットは41種類の薬剤，元画像合計287枚である．
薬剤ごとの内訳を表\ref{tab:dataset}に示す．

\begin{table}[htbp]
\centering
\caption{データセットの内訳（41薬剤）．学習・評価枚数は11種データ拡張適用後の枚数}
\label{tab:dataset}
\scriptsize
\setlength{\tabcolsep}{3pt}
\begin{tabular}{lrrr|lrrr}
\hline
薬剤名 & 元 & 学習 & 評価 & 薬剤名 & 元 & 学習 & 評価 \\
\hline
メインテート錠2.5mg & 8 & 66 & 22 & トラゼンタ錠5mg & 7 & 55 & 22 \\
アルダクトンＡ錠25mg & 6 & 44 & 22 & ノルバスクOD錠5mg & 6 & 44 & 22 \\
ウルソ錠100mg & 4 & 33 & 11 & バイアスピリン錠100mg & 4 & 33 & 11 \\
エソメプラゾールカプセル20mg & 7 & 55 & 22 & バクトラミン配合錠 & 6 & 44 & 22 \\
エンレスト錠100mg & 9 & 77 & 22 & パルモディアXR錠0.2mg & 10 & 88 & 22 \\
エンレスト錠200mg & 9 & 77 & 22 & フェブリク錠10mg & 9 & 77 & 22 \\
カロナール錠500mg & 6 & 44 & 22 & フォシーガ錠10mg & 9 & 77 & 22 \\
キャブピリン配合錠 & 9 & 77 & 22 & フロセミド錠20mg & 6 & 44 & 22 \\
クエン酸第一鉄Na錠50mg & 6 & 44 & 22 & プレドニン錠5mg & 6 & 44 & 22 \\
ジャディアンス錠10mg & 6 & 44 & 22 & プロマックD錠75mg & 9 & 77 & 22 \\
ジャヌビア錠50mg & 9 & 77 & 22 & ベオーバ錠50mg & 9 & 77 & 22 \\
センノシド錠12mg & 9 & 77 & 22 & ミヤＢＭ錠 & 4 & 33 & 11 \\
タケキャブ錠10mg & 7 & 55 & 22 & ムコダイン錠500mg & 9 & 77 & 22 \\
タケキャブ錠20mg & 4 & 33 & 11 & メチコバール錠500$\mu$g & 4 & 33 & 11 \\
タリージェOD錠5mg & 6 & 44 & 22 & メトクロプラミド錠5mg & 4 & 33 & 11 \\
ダイアート錠30mg & 10 & 88 & 22 & メトグルコ錠250mg & 10 & 88 & 22 \\
デエビゴ錠5mg & 10 & 88 & 22 & ランソプラゾールOD錠15mg & 4 & 33 & 11 \\
リクシアナOD錠30mg & 9 & 77 & 22 & レバミピド錠100mg & 6 & 44 & 22 \\
レボチロキシンNa錠50$\mu$g & 10 & 88 & 22 & レボフロキサシン錠500mg & 3 & 22 & 11 \\
ロスバスタチン錠2.5mg「DSEP」 & 6 & 44 & 22 & ロキソプロフェンナトリウム錠60mg & 6 & 44 & 22 \\
酸化マグネシウム錠330mg & 6 & 44 & 22 & \textbf{合計} & \textbf{287} & \textbf{2,343} & \textbf{814} \\
\hline
\end{tabular}
\end{table}

\subsubsection{学習・評価分割}

学習データと評価データの分離は，データ拡張を行う前の元画像単位で実施した．
各薬剤フォルダごとに，画像ファイル名を辞書順に整列したうえで乱数シード42でシャッフルし，先頭75\%を学習用，残り25\%を評価用とした（評価用には最低1枚を保証）．
シードを固定しているため，同一の元画像集合からは常に同一の分割が再現される．
分割を元画像単位で行うため，同一の元画像に由来する拡張画像が学習データと評価データの両方に含まれることはない．

\subsubsection{データ拡張}

撮影枚数が限定的であることを補うために，分割後の学習用・評価用それぞれの元画像に対して，表\ref{tab:aug}に示す11種類のデータ拡張を適用し，PNG形式で保存した．

\begin{table}[htbp]
\centering
\caption{データ拡張の種類とパラメータ（1枚の元画像から11枚を生成）}
\label{tab:aug}
\small
\begin{tabular}{llp{0.45\linewidth}}
\hline
拡張名 & 種別 & 処理内容 \\
\hline
orig & 複製 & 元画像のコピー \\
rot$\pm$15 & 回転 & $\pm$15度回転（キャンバス拡張あり，余白はRGB(128,128,128)で充填） \\
rot$\pm$30 & 回転 & $\pm$30度回転（同上） \\
flip\_h & 反転 & 水平（左右）反転 \\
bright0.7 / 1.3 & 明度 & 明度を0.7倍／1.3倍に変更 \\
contrast0.7 / 1.3 & コントラスト & コントラストを0.7倍／1.3倍に変更 \\
noise & ノイズ & 平均0・標準偏差10のガウシアンノイズを加算（画素値は0〜255にクリップ） \\
\hline
\end{tabular}
\end{table}

この結果，学習用2,343枚（元213枚$\times$11），評価用814枚（元74枚$\times$11）のデータセットを構築した．

\subsection{学習設定}

基盤モデルとしてQwen3.5-4Bを用い，\ref{sec3}.4節で述べたLoRAファインチューニングを行った．
学習に用いたハイパーパラメータを表\ref{tab:lora_config}に示す．

\begin{table}[htbp]
\centering
\caption{LoRAファインチューニングの設定}
\label{tab:lora_config}
\begin{tabular}{ll}
\hline
パラメータ & 値 \\
\hline
基盤モデル & Qwen/Qwen3.5-4B \\
LoRA rank $r$ & 8 \\
LoRA alpha $\alpha$ & 16 \\
target modules & q\_proj，v\_proj \\
LoRA dropout & 0.05 \\
エポック数 & 5 \\
最適化手法 & AdamW \\
学習率 & 2e-4（固定） \\
バッチサイズ & 1 \\
勾配チェックポイント & 有効 \\
演算精度 & bfloat16 \\
画像入力ピクセル数 & 最小 $256 \times 256$，最大 $768 \times 768$ \\
\hline
\end{tabular}
\end{table}

学習サンプルは「ユーザ発話＝薬剤画像＋質問文，アシスタント応答＝正解薬剤名（フォルダ名）＋文末トークン」のチャット形式で構成し，\ref{sec3}.4節で述べたとおり質問プロンプト部分は学習の対象とせず，応答部分（正解薬剤名）のみを学習させた．
各エポックで学習サンプルの順序をシャッフルした．
41薬剤・2,343枚での学習に要した時間は498分50秒（約8.3時間）であり，エポックごとの平均損失は表\ref{tab:loss_all}のとおり単調に減少した．

\begin{table}[htbp]
\centering
\caption{全薬剤学習における損失推移（学習2,343枚）}
\label{tab:loss_all}
\begin{tabular}{rr}
\hline
エポック & 平均損失 \\
\hline
1 & 0.2470 \\
2 & 0.0258 \\
3 & 0.0133 \\
4 & 0.0058 \\
5 & 0.0006 \\
\hline
\end{tabular}
\end{table}

\subsection{評価手順}\label{sec:evalproc}

評価時には学習時と同一の質問文「この錠剤は何という薬ですか？薬剤名だけで答えてください。」でモデルに推論を実行させ，表記のゆれ（全角・半角，大文字・小文字）を揃えたうえで，正解薬剤名（フォルダ名）が回答テキストに含まれていれば正解と判定した．
Top-3評価では，ビームサーチ（ビーム幅3）により3候補を生成させた．

比較手法の評価条件は次のとおりである．
ファインチューニング前のQwen3.5-4Bは，LoRAを適用しない同一の基盤モデルに同一の質問文・同一の評価画像で推論させた．
Geminiは\texttt{gemini-3.5-flash}をAPI経由で用い，プロンプトには「PTPシートに印刷された文字を読み，製品名と用量（mg/$\mu$g）を必ず両方答えてください」を使用した．
APIのレート制限対策として1リクエストごとに1秒の待機を挟み，レート超過エラー（429）発生時には再試行した．
正解判定はQwen側と同一の正規化・部分一致判定である．

\subsection{予備実験（単一薬剤メインテート錠）}

本実験に先立ち，単一薬剤「メインテート錠2.5mg」のみを対象として，（1）出力形式・方式の比較，（2）少数画像＋データ拡張でのLoRA学習の有効性確認を行った．

\subsubsection{出力形式・方式の比較}

提案手法の出力形式を確定するため，以下の4方式をテスト画像5枚（メインテート錠2.5mg）で比較した（2026年4月26日実施）．

\begin{itemize}
  \item \textbf{方法A（薬剤名のみFT）}：画像に対して薬剤名のみを出力するようLoRA学習する方式（\ref{sec3}.4節の提案手法）．
  \item \textbf{方法B（詳細鑑別FT）}：薬剤名に加えて色・形状・刻印などの鑑別理由を含む詳細な説明を出力するようLoRA学習する方式．
  \item \textbf{RAG型}：薬剤情報データベース（JSON形式の薬剤DB）を検索し，候補薬剤情報を回答に利用する方式．
  \item \textbf{Gemini}：外部APIによるゼロショット識別（当時は\texttt{gemini-2.5-flash}を使用）．
\end{itemize}

結果を表\ref{tab:method_compare}に示す．

\begin{table}[htbp]
\centering
\caption{出力形式・方式の比較（テスト画像5枚・メインテート錠2.5mg）}
\label{tab:method_compare}
\begin{tabular}{lrrl}
\hline
方式 & 成功数 & 平均時間/枚 & 備考 \\
\hline
\textbf{方法A（薬剤名のみFT）} & \textbf{5 / 5} & 0.8 秒 & 最も高速かつ正確 \\
方法B（詳細鑑別FT） & 3 / 5 & 37.4 秒 & 説明的だが遅く不安定 \\
RAG型（薬剤DB検索） & 5 / 5 & 24.4 秒 & 正確だが検索処理で低速 \\
Gemini（ゼロショット） & 0 / 5 & 3.4 秒 & 全滅 \\
\hline
\end{tabular}
\end{table}

方法AとRAG型がともに5/5で正答したが，推論時間は方法Aが0.8秒/枚と約30倍高速であったため，以降の全実験では方法A（薬剤名のみを出力）を採用した．
なお，RAG型は新規薬剤の追加に再学習が不要という利点があり，全薬剤規模での比較は今後の課題とする（\ref{sec5}章）．

\subsubsection{拡張データセットでの学習と評価}

次に，少数の元画像とデータ拡張の組み合わせでLoRA学習が機能するかを確認した．
元画像6枚を11種拡張した66枚で学習し，元画像2枚を11種拡張した22枚で評価した．
学習に要した時間は15分18秒であり，エポックごとの平均損失は表\ref{tab:loss_meinte}のとおりであった．

\begin{table}[htbp]
\centering
\caption{単一薬剤での学習損失推移}
\label{tab:loss_meinte}
\begin{tabular}{rr}
\hline
エポック & 平均損失 \\
\hline
1 & 0.6757 \\
2 & 0.0002 \\
3 & 0.0001 \\
4 & 0.0001 \\
5 & 0.0001 \\
\hline
\end{tabular}
\end{table}

2エポック目以降にほぼ完全に収束し，評価用22枚すべてで正解薬剤名を出力した（Top-1精度100\%）．
明度，コントラスト，回転，水平反転，ノイズの全拡張パターンに対して正答した．
同一の22枚に対するGemini（\texttt{gemini-2.5-flash}）のゼロショット識別は0/22であった（表\ref{tab:vs_gemini}）．
ただしGemini側の失敗には無料枠のAPIリクエスト数上限超過によるエラー応答6件が含まれるため，この予備実験の比較は参考値であり，本実験（\ref{sec:single_result}節）にて全薬剤・新プロンプトでの再評価を行った．

\begin{table}[htbp]
\centering
\caption{メインテート錠2.5mgにおけるLoRAとGeminiの比較（予備実験・参考値）}
\label{tab:vs_gemini}
\begin{tabular}{lrr}
\hline
指標 & LoRA & Gemini 2.5 Flash \\
\hline
成功数 & 22 / 22 & 0 / 22 \\
成功率 & 100.0\% & 0.0\% \\
平均推論時間 & 1.1 秒 & 14.0 秒 \\
\hline
\end{tabular}
\end{table}

\subsection{初期検討：データ拡張なしでの全薬剤学習}\label{sec:noaug}

データ拡張を導入する前の初期検討として，元画像287枚のみ（拡張なし）を用いた全薬剤のLoRA学習を行った（2026年5月12日頃実施）．
学習設定は本実験と同一（LoRA rank 8・alpha 16，target modules q\_proj・v\_proj，5エポック，学習率2e-4）である．

学習損失は初期3.88から最終約0.001までほぼゼロに低下した．
学習枚数が287枚と少ないにもかかわらず損失がほぼゼロまで落ちる挙動は，学習データを丸暗記する過学習の兆候であり，未学習画像への汎化は期待できないと判断した．
このモデルに対する定量評価は実施せず，撮影枚数の不足を補う対策として\ref{sec:dataset}節の11種類のデータ拡張を導入した．
すなわち本検討は，提案手法においてデータ拡張が必要である根拠を与えた通過点である．
以降の結果はすべて拡張あり（学習2,343枚）のモデルによる．

\subsection{単一薬剤識別の結果（本実験）}\label{sec:single_result}

41薬剤の評価画像814枚（撮影＋拡張あり）に対して，提案手法（ファインチューニング後のQwen3.5-4B），ファインチューニング前のQwen3.5-4B，およびGeminiの3者を同一条件で比較した．
結果を表\ref{tab:single_compare}に示す．

\begin{table}[htbp]
\centering
\caption{単一薬剤識別の3者比較（同一の評価画像814枚）}
\label{tab:single_compare}
\begin{tabular}{lrr}
\hline
手法 & Top-1精度 & 成功数 \\
\hline
\textbf{ファインチューニング後 Qwen3.5-4B（提案手法）} & \textbf{98.2\%} & 799 / 814 \\
Gemini（\texttt{gemini-3.5-flash}・ゼロショット） & 82.3\% & 670 / 814 \\
ファインチューニング前 Qwen3.5-4B（ゼロショット） & 0.9\% & 7 / 814 \\
\hline
\end{tabular}
\end{table}

提案手法は98.2\%を達成し，41薬剤中29薬剤で100\%の正答であった．
提案手法の薬剤別の成功率を表\ref{tab:per_drug}に示す．

\begin{table}[htbp]
\centering
\caption{提案手法（ファインチューニング後4B）の薬剤別Top-1成功率（成功数／評価枚数）}
\label{tab:per_drug}
\scriptsize
\setlength{\tabcolsep}{3pt}
\begin{tabular}{lrr|lrr}
\hline
薬剤名 & 成功／枚数 & 成功率 & 薬剤名 & 成功／枚数 & 成功率 \\
\hline
メインテート錠2.5mg & 22/22 & 100.0\% & パルモディアXR錠0.2mg & 20/22 & 90.9\% \\
アルダクトンＡ錠25mg & 21/22 & 95.5\% & フェブリク錠10mg & 22/22 & 100.0\% \\
ウルソ錠100mg & 11/11 & 100.0\% & フォシーガ錠10mg & 22/22 & 100.0\% \\
エソメプラゾールカプセル20mg & 21/22 & 95.5\% & フロセミド錠20mg & 21/22 & 95.5\% \\
エンレスト錠100mg & 21/22 & 95.5\% & プレドニン錠5mg & 22/22 & 100.0\% \\
エンレスト錠200mg & 22/22 & 100.0\% & プロマックD錠75mg & 22/22 & 100.0\% \\
カロナール錠500mg & 22/22 & 100.0\% & ベオーバ錠50mg & 22/22 & 100.0\% \\
キャブピリン配合錠 & 22/22 & 100.0\% & ミヤＢＭ錠 & 11/11 & 100.0\% \\
クエン酸第一鉄Na錠50mg & 22/22 & 100.0\% & ムコダイン錠500mg & 21/22 & 95.5\% \\
ジャディアンス錠10mg & 22/22 & 100.0\% & メチコバール錠500$\mu$g & 11/11 & 100.0\% \\
ジャヌビア錠50mg & 20/22 & 90.9\% & メトクロプラミド錠5mg & 11/11 & 100.0\% \\
センノシド錠12mg & 22/22 & 100.0\% & メトグルコ錠250mg & 22/22 & 100.0\% \\
タケキャブ錠10mg & 22/22 & 100.0\% & ランソプラゾールOD錠15mg & 11/11 & 100.0\% \\
タケキャブ錠20mg & 11/11 & 100.0\% & リクシアナOD錠30mg & 21/22 & 95.5\% \\
タリージェOD錠5mg & 21/22 & 95.5\% & レバミピド錠100mg & 22/22 & 100.0\% \\
ダイアート錠30mg & 22/22 & 100.0\% & レボチロキシンNa錠50$\mu$g & 22/22 & 100.0\% \\
デエビゴ錠5mg & 22/22 & 100.0\% & レボフロキサシン錠500mg & 10/11 & 90.9\% \\
トラゼンタ錠5mg & 22/22 & 100.0\% & ロキソプロフェンナトリウム錠60mg & 22/22 & 100.0\% \\
ノルバスクOD錠5mg & 20/22 & 90.9\% & ロスバスタチン錠2.5mg「DSEP」 & 21/22 & 95.5\% \\
バイアスピリン錠100mg & 11/11 & 100.0\% & 酸化マグネシウム錠330mg & 22/22 & 100.0\% \\
バクトラミン配合錠 & 22/22 & 100.0\% & \textbf{全体} & \textbf{799/814} & \textbf{98.2\%} \\
\hline
\end{tabular}
\end{table}

ファインチューニング前のモデルは0.9\%（7/814）であり，正答した7枚はすべてレバミピド錠100mgであった（同薬剤のみ31.8\%，他の40薬剤は0\%）．
レバミピドのみ正答できたのは，一般名としての知名度が高く事前学習の知識に含まれていたためと推測される．
ファインチューニング前後の差97.3ポイントが，薬剤画像特化の追加学習の効果である．

Geminiは82.3\%であり，提案手法が15.9ポイント上回った．
Geminiは薬剤ごとの差が大きく，提案手法では高精度に識別できる薬剤で大きく精度を落とした（表\ref{tab:gemini_weak}）．

\begin{table}[htbp]
\centering
\caption{Geminiの正答率が低かった薬剤と提案手法との比較（同一評価画像）}
\label{tab:gemini_weak}
\small
\begin{tabular}{lrr}
\hline
薬剤名 & Gemini & 提案手法（FT後4B） \\
\hline
パルモディアXR錠0.2mg & 4.5\% & 90.9\% \\
ジャヌビア錠50mg & 18.2\% & 90.9\% \\
メチコバール錠500$\mu$g & 36.4\% & 100.0\% \\
エンレスト錠200mg & 45.5\% & 100.0\% \\
カロナール錠500mg & 45.5\% & 100.0\% \\
ロキソプロフェンナトリウム錠60mg & 50.0\% & 100.0\% \\
\hline
\end{tabular}
\end{table}

なお，Geminiの評価は当初「薬剤名のみを日本語で答える」旧プロンプトで実施しており，このときの正答率は44.0\%（学習用拡張データ1,338枚を対象とした途中中断時点の参考値）であった．
プロンプトを「PTPシートの印字を読み，製品名と用量を必ず両方答える」形式（\ref{sec:evalproc}節）に変更して評価用814枚で測り直した結果が82.3\%である．
外部APIの性能はプロンプト設計に強く依存することも確認された．

また，提案手法について評価条件を変えた結果を表\ref{tab:single_extra}に示す．

\begin{table}[htbp]
\centering
\caption{提案手法（ファインチューニング後4B）の評価条件別の結果}
\label{tab:single_extra}
\begin{tabular}{lrr}
\hline
評価条件 & 評価画像 & 精度 \\
\hline
撮影＋拡張あり（標準） & 814枚 & Top-1 98.2\% \\
拡張なし（元画像のみ） & 74枚 & Top-1 100.0\% \\
Top-3評価（ビーム幅3） & 814枚 & Top-3 99.9\%（Top-1 98.3\%） \\
\hline
\end{tabular}
\end{table}

拡張なしの元画像のみでは100\%（74/74）であり，98.2\%における失敗15枚はすべて拡張によって人工的に変形させた画像に由来する．
また，上位3候補を提示するTop-3評価では99.9\%（813/814）に達しており，Top-1で誤る場合でも候補提示により実用上の救済が可能である．

誤識別15枚は表\ref{tab:per_drug}に示した12薬剤で発生した（各薬剤1〜2枚）．
この15枚を拡張種別ごとに集計したものを表\ref{tab:failed_aug}に示す．

\begin{table}[htbp]
\centering
\caption{誤識別画像の拡張種別内訳（提案手法・15枚）}
\label{tab:failed_aug}
\begin{tabular}{lr}
\hline
拡張種別 & 失敗枚数 \\
\hline
回転 +30度 & 4 \\
回転 -30度 & 4 \\
水平反転 & 4 \\
回転 -15度 & 2 \\
明度 1.3倍 & 1 \\
元画像，回転 +15度，明度 0.7倍，コントラスト 0.7・1.3倍，ノイズ & 0 \\
\hline
合計 & 15 \\
\hline
\end{tabular}
\end{table}

誤識別の大部分は回転（$\pm$30度，$-$15度）と水平反転の合計14枚であり，文字方向や錠剤シートの上下左右が大きく変化する拡張に対してモデルが脆弱であることが分かった．
一方で，明度・コントラスト・ノイズなどの輝度系の変化に対する失敗は1件のみであり，色や明るさの揺らぎに対しては頑健であった．
また，同一の元画像が複数の拡張で同時に失敗するケースが3例観察され（ジャヌビア錠50mg，ノルバスクOD錠5mg，パルモディアXR錠0.2mgの各1枚），これらはそもそも判別が難しい撮影条件の画像であり，撮影品質に起因する可能性がある．
誤識別の出力内容は無関係な単語ではなくほぼすべて別の薬剤名であり，回転や反転で文字が読みにくくなった結果，近い見た目の別薬剤に誤判定していると考えられる．

\subsection{外部画像での汎化（予備的検証）}

評価画像が学習画像と同一の撮影環境であることによる過大評価の可能性を確認するため，データセット（\ref{sec:dataset}節）とは別の時期に撮影したレバミピド錠100mgの画像を，回転・明度の拡張を施して390枚とし，提案手法（同一の学習済みLoRA・同一の評価手順）で評価した（2026年5月21日実施）．
結果は99.7\%（389/390，失敗1枚）であり，少なくとも本薬剤については撮影時期が異なっても高精度を維持した．
ただし対象が1薬剤のみであるため，全薬剤を対象とした別環境画像（Web収集画像等）での汎化評価は今後の課題である（\ref{sec5}章）．

\subsection{複数薬剤同時識別の結果}\label{sec:multi_result}

\subsubsection{評価データの作成}

1枚に複数の薬剤が写る評価画像は，評価用814枚（学習に未使用）の画像から合成して作成した．
合成手順は次のとおりである．
（1）各薬剤画像から，画素値の局所標準偏差に基づくテクスチャ検出により薬剤（PTPシート）部分の領域を自動的に切り出す．
（2）切り出した薬剤画像を高さ512pxに揃え，異なる薬剤2枚または3枚を余白なしで横に連結して1枚のタイル画像とする．
（3）全41薬剤が均等に登場するよう，各画像の主薬剤をラウンドロビンで割り当て，残りの1〜2薬剤をランダムに選択する（乱数シード42で固定）．
この方法で2薬剤500枚＋3薬剤500枚の合計1,000枚を生成した（各薬剤は各グループに12〜13回ずつ登場）．
各合成画像に写る薬剤の集合は正解ラベルファイル（\texttt{labels.json}）として保存し，採点に用いた．
合成には評価用画像のみを使用しているため，学習データのリークはない．

\subsubsection{一括質問方式の失敗}

まず，合成画像全体をモデルに1回提示し「写っている薬をすべてJSONリストで答えよ」と質問する一括質問方式を試みた．
表\ref{tab:multi_naive}に示すとおり，プロンプトの工夫（必ず複数答えるよう指示，薬剤の個数を提示，領域順の列挙指示，高解像度化）のいずれを加えても完全一致率はほぼ0\%であった．

\begin{table}[htbp]
\centering
\caption{一括質問方式の結果（ファインチューニング後4B）}
\label{tab:multi_naive}
\small
\begin{tabular}{lp{0.4\linewidth}r}
\hline
方式 & プロンプトの工夫 & 完全一致率 \\
\hline
手法1 & 基準（すべてJSONで答えよ） & $\approx$0\%（1/2050） \\
手法3 & 「必ず複数答えよ」と指示 & 0\% \\
手法4 & 写っている個数$n$を提示 & 0\% \\
手法5 & 「領域ごとに左から順に」と指示 & 0\% \\
手法6 & 個数提示＋高解像度入力 & 0\% \\
\hline
\end{tabular}
\end{table}

失敗の原因は，モデルが「1枚に1薬剤」を答えるよう学習されているため，複数の薬剤が写っていても1薬剤しか出力しないことであった．

\subsubsection{検出→単剤分類方式の結果と3者比較}

次に，\ref{sec:multi}節の2段構成（検出→単剤分類）を適用した．
検出にはLoRAを無効化した素のベースモデルのゼロショット検出（検出プロンプトは\ref{sec:multi}節，入力最大$1280 \times 1280$ピクセル）を用い，検出された各領域をマージン2\%付きで切り出して単剤識別した．

まず先行検証として，学習用画像から同じ手順で合成した生成画像2,050枚（各薬剤が2薬剤25枚＋3薬剤25枚の計50回登場）で本方式を評価し，完全一致率69.6\%・F1 85.2\%を得た．
ただしこの2,050枚は学習に使用した画像を合成元とするため評価としては甘い．
そこで本評価では，学習に未使用の評価用画像から合成した前述の1,000枚を用いた．
同一の1,000枚で，Geminiおよびファインチューニング前のQwen3.5-4Bとも比較した．
Geminiは外部APIであり検出→分類のパイプラインを組めないため，一括質問方式（プロンプト「この画像に写っているすべての薬を薬剤名と用量を含めてJSONリストで答えてください」）で評価した．
ファインチューニング前のQwen3.5-4Bも同様に一括質問方式（「この画像に写っているすべての薬をJSONリストで答えてください」）である．
結果を表\ref{tab:multi_compare}に示す．

\begin{table}[htbp]
\centering
\caption{複数薬剤同時識別の3者比較（同一の合成評価画像1,000枚）}
\label{tab:multi_compare}
\small
\begin{tabular}{lrrrr}
\hline
手法 & 完全一致率 & Recall & Precision & F1 \\
\hline
\textbf{FT後4B（検出→単剤分類・提案）} & \textbf{68.8\%}（688/1000） & 85.5\% & 86.3\% & 85.6\% \\
Gemini（一括質問） & 50.6\%（506/1000） & --- & --- & --- \\
FT前4B（一括質問） & 0.0\%（0/1000） & --- & --- & --- \\
\hline
\end{tabular}
\end{table}

提案手法の完全一致率は68.8\%（2薬剤75.6\%，3薬剤62.0\%）であり，Geminiの50.6\%（2薬剤60.4\%，3薬剤40.8\%）を上回った．
ファインチューニング前のモデルは完全一致0\%（1薬剤以上正解した画像でも3.3\%）であった．
単剤識別が98.2\%であるのに対し複数同時識別が68.8\%に留まることから，性能低下の主因は検出・切り出し段階にあると考えられた．
この仮説は次節の実験で検証する．

\subsection{検出器の専用学習による改善}\label{sec:detector}

\subsubsection{評価セットと上限値の診断（検出が完璧な場合）}

本節の実験では，評価用814枚（学習に未使用）から各薬剤の正解bounding box付きで再合成した評価セットを用いた．
構成は，標準セット200枚（2薬剤100枚＋3薬剤100枚）と，評価用814枚を1枚ずつ使い切る全部セット320枚（2薬剤146枚＋3薬剤174枚）の2種類である．

検出段階の改善余地を定量化するため，合成時に保存した正解bounding boxで切り出して単剤識別する「検出が完璧な場合の上限値」を測定した．
あわせて合成タイル画像の解像度（タイル高さ）の影響も調べた．
結果を表\ref{tab:a0}に示す．

\begin{table}[htbp]
\centering
\caption{合成解像度別の上限値（正解bbox使用）と実機値（素のベース検出）}
\label{tab:a0}
\begin{tabular}{lrr}
\hline
タイル高さ & 上限（正解bbox） & 実機（素のベース検出） \\
\hline
512px & 80.0\% & --- \\
768px & 92.0\% & 81.5\%（200枚） \\
1024px & 93.5\% & --- \\
\hline
\end{tabular}
\end{table}

解像度を512pxから768pxに上げるだけで上限が80.0\%から92.0\%に，実機（素のベース検出）も68.8\%から81.5\%に改善した．
したがって（1）入力解像度と（2）検出精度の双方が複数識別のボトルネックであることが確認された．

\subsubsection{検出器LoRAの学習}

検出精度を高めるため，検出専用のLoRAアダプタを追加学習した．
学習データは，学習用分割の画像（評価用合成画像とは元画像レベルで分離しリークを防止）から，実写に近い条件（薬剤間のすきま，背景余白約6\%，大きさ・位置のばらつき）を加えて合成した複数薬剤画像500枚（2薬剤200枚・3薬剤200枚・4薬剤100枚）である．
教師信号は，各薬剤のbounding boxを0〜1000に正規化した座標で記述するグラウンディング形式JSON（\texttt{[\{"bbox\_2d":[x1,y1,x2,y2],"label":"薬剤"\}]}）であり，検出プロンプトは推論時（\ref{sec:multi}節）と同一とした．
学習設定はLoRA rank 8・alpha 16，画像解像度1024px，3エポックで，平均損失は0.181 → 0.134 → 0.119と単調に減少した（所要約80分）．

\subsubsection{検出器差し替え後の再評価}

学習した検出器LoRAを検出段に差し替え，同一の合成画像で「上限（正解bbox）」「旧検出（素のベース）」「新検出（学習した検出器）」を比較した．
1つの基盤モデルに識別LoRAと検出LoRAの2つを搭載し，処理段階に応じてアダプタを切り替えた．
結果を表\ref{tab:a3}に示す．

\begin{table}[htbp]
\centering
\caption{検出器学習の効果（完全一致率，括弧内はF1）}
\label{tab:a3}
\small
\setlength{\tabcolsep}{4pt}
\begin{tabular}{lrrr}
\hline
方式 & 768px・標準200枚 & 768px・全部320枚 & 1024px・標準200枚 \\
\hline
上限（正解bbox） & 92.0\%（96.8\%） & 91.6\%（96.9\%） & 93.5\%（97.3\%） \\
旧検出（素のベース） & 81.5\%（90.8\%） & 77.8\%（90.4\%） & 85.5\%（93.4\%） \\
\textbf{新検出（学習した検出器）} & \textbf{91.0\%}（96.2\%） & \textbf{88.8\%}（96.3\%） & \textbf{96.0\%}（98.4\%） \\
改善（新$-$旧） & +9.5pt & +10.9pt & +10.5pt \\
\hline
\end{tabular}
\end{table}

検出器の専用学習により完全一致率は9.5〜10.9ポイント改善し，検出起因の損失の約80〜90\%を回収した．
新検出（学習した検出器）の薬剤数別の内訳を表\ref{tab:a3_group}に示す．

\begin{table}[htbp]
\centering
\caption{新検出（学習した検出器）の薬剤数別の完全一致率}
\label{tab:a3_group}
\small
\begin{tabular}{lrr}
\hline
条件 & 2薬剤 & 3薬剤 \\
\hline
768px・標準200枚 & 92.0\%（92/100） & 90.0\%（90/100） \\
768px・全部320枚 & 93.8\%（137/146） & 84.5\%（147/174） \\
1024px・標準200枚 & 97.0\%（97/100） & 95.0\%（95/100） \\
\hline
\end{tabular}
\end{table}

特に1024px設定では全体96.0\%（2薬剤97.0\%，3薬剤95.0\%）に達し，本研究の複数薬剤同時識別における最高値となった．
注目すべきは，1024pxにおいて新検出（96.0\%）が正解bbox使用の上限（93.5\%）を上回った点である．
これは，学習した検出器が出力するbounding boxが正解box（タイル全体）よりも薬剤に密着しており，切り出し品質が高いためである．
1024pxで新検出が誤った8枚はすべて検出ではなく識別段階の取り違え（うち5枚がロスバスタチン錠2.5mg「DSEP」→クエン酸第一鉄Na錠50mgの誤読）であり，検出の問題は実質的に解消された．

\subsection{過検出抑制の検証（負の結果）}

検出器学習前の誤り分析で観察された過検出（余分なboxの出力）への対策として，（1）重なり率（IoU）0.5超のboxを統合する重複除去（NMS），（2）切り出し領域が薬剤か否かを判定する「該当なしゲート」LoRA（薬剤328枚・非薬剤340枚の計668枚で学習）の2方式を，768px・標準200枚で検証した．
結果を表\ref{tab:a4}に示す．

\begin{table}[htbp]
\centering
\caption{過検出抑制の検証結果（768px・標準200枚）}
\label{tab:a4}
\begin{tabular}{lrrr}
\hline
方式 & 完全一致率 & F1 & Precision \\
\hline
ベース（A-3新検出・抑制なし） & 91.0\% & 96.2\% & 96.5\% \\
重複除去（NMS） & 91.0\% & 96.2\% & 96.5\% \\
該当なしゲート & 83.5\% & 94.3\% & 96.5\% \\
\hline
\end{tabular}
\end{table}

重複除去は除去対象となる重複boxが0個であり効果がなかった．
該当なしゲートは学習損失が2エポック目で0.0000まで低下する過学習を起こし，実際の検出領域（特にジャディアンス錠10mg）を誤って「該当なし」と判定して正しい薬剤を棄却した結果，完全一致率が7.5ポイント悪化した（Precisionは向上せず）．
この負の結果から，過検出は検出器の専用学習の時点で既に解消されており，後処理による抑制の余地はないことが確認された．

\subsection{軽量化の評価（量子化・モデルサイズ）}\label{sec:light}

スマートフォン等の端末搭載を見据え，（1）学習済み4Bモデルの4bit量子化（NF4形式＋二重量子化．再学習は行わず，学習済みLoRAをそのまま量子化モデルに適用），（2）モデルサイズの縮小（0.8B / 2B，いずれも4Bと同一の学習データ・同一のLoRA設定で学習）による精度の変化を測定した．

評価は3条件で行った．
単一薬剤は評価用814枚（\ref{sec:single_result}節と同一），複数薬剤は合成評価画像1,000枚（\ref{sec:multi_result}節と同一・検出→単剤分類方式）である．
混在評価は，薬剤数が事前に分からない実運用を模擬するため，単一薬剤の評価画像から100枚と複数薬剤の合成画像から100枚を抽出・シャッフルした計200枚に対し，検出→単剤分類方式で薬剤数を自動判定させて評価した．
結果を表\ref{tab:size_single}〜\ref{tab:size_mixed}に示す．

\begin{table}[htbp]
\centering
\caption{モデルサイズ・量子化と単一薬剤識別精度（Top-1，814枚）}
\label{tab:size_single}
\begin{tabular}{lrrr}
\hline
モデル & 非量子化 & 4bit量子化 & 量子化での低下 \\
\hline
0.8B & 87.3\%（711/814） & 74.3\%（605/814） & $-$13.0pt \\
2B & 87.2\%（710/814） & 83.9\%（683/814） & $-$3.3pt \\
\textbf{4B} & \textbf{98.2\%}（799/814） & \textbf{91.9\%}（748/814） & $-$6.3pt \\
\hline
\end{tabular}
\end{table}

\begin{table}[htbp]
\centering
\caption{モデルサイズ・量子化と複数薬剤同時識別精度（完全一致率，括弧内はF1，1,000枚）}
\label{tab:size_multi}
\begin{tabular}{lrr}
\hline
モデル & 非量子化 & 4bit量子化 \\
\hline
0.8B & 12.7\%（54.1\%） & 4.6\%（33.9\%） \\
2B & 30.0\%（67.6\%） & 12.9\%（51.7\%） \\
\textbf{4B} & \textbf{68.8\%}（85.6\%） & \textbf{55.2\%}（76.7\%） \\
\hline
\end{tabular}
\end{table}

\begin{table}[htbp]
\centering
\caption{単一・複数混在評価（完全一致率，括弧内はF1，単一100枚＋複数100枚）}
\label{tab:size_mixed}
\small
\begin{tabular}{lrr}
\hline
モデル & 非量子化（単一／複数の内訳） & 4bit量子化（単一／複数の内訳） \\
\hline
0.8B & 38.0\%・F1 62.7\%（66\%／10\%） & 26.5\%・F1 42.4\%（48\%／5\%） \\
2B & 48.5\%・F1 68.3\%（68\%／29\%） & 40.5\%・F1 62.2\%（71\%／10\%） \\
\textbf{4B} & \textbf{68.0\%}・F1 80.3\%（67\%／69\%） & \textbf{70.0\%}・F1 80.4\%（72\%／68\%） \\
\hline
\end{tabular}
\end{table}

単一薬剤識別ではモデルサイズによる差が比較的小さい（0.8Bでも87.3\%）のに対し，複数薬剤同時識別では0.8B 12.7\%，2B 30.0\%，4B 68.8\%とサイズの影響が極めて大きい．
量子化の影響も複数識別の方が大きく（4Bで単一$-$6.3pt，複数$-$13.6pt），これは検出→分類の2段処理で量子化誤差が領域ごとに累積するためと考えられる．
また，量子化による精度低下は薬剤間で一様ではなく，量子化4Bの単一識別ではレバミピド錠100mg（27.3\%），メチコバール錠500$\mu$g・ロスバスタチン錠2.5mg「DSEP」（各54.5\%）など，特定の薬剤に低下が集中した．
量子化4Bの複数識別の内訳は2薬剤63.8\%・3薬剤46.6\%（Recall 76.8\%・Precision 77.6\%）であった．
混在評価では4Bのみが単一・複数で均等な精度を示し（67\%／69\%），量子化しても全体精度がほぼ低下しなかった（68.0\%→70.0\%，各100枚での誤差範囲）．
量子化4Bでも複数識別55.2\%はGeminiの50.6\%を上回っており，端末搭載の候補としては量子化4Bが精度と軽さのバランスで最有力である．
なお混在評価は各100枚の軽量版であり，数ポイントの差は誤差として傾向の確認に用いる．

\subsection{誤識別の横断分析}\label{sec:error}

単一薬剤識別の全6条件（モデルサイズ3種$\times$量子化の有無）の誤識別を横断的に集計し，誤りの傾向を分析した．
条件別の誤識別件数を表\ref{tab:err_by_model}に示す（合計640件）．

\begin{table}[htbp]
\centering
\caption{単一薬剤識別の条件別誤識別件数（評価814枚に対する誤り数）}
\label{tab:err_by_model}
\begin{tabular}{lrr}
\hline
モデル & 非量子化 & 4bit量子化 \\
\hline
0.8B & 103 件 & 209 件 \\
2B & 104 件 & 131 件 \\
4B & 27 件 & 66 件 \\
\hline
\end{tabular}
\end{table}

\textbf{（1）向きの崩れへの脆弱性}：
誤識別を評価画像の拡張種別ごとに集計した結果を表\ref{tab:err_by_aug}に示す．
回転（$\pm$15度・$\pm$30度）合計311件と水平反転108件を合わせた419件（全体の約66\%）が，向きが崩れる加工の画像で発生した．
明度・コントラスト・ノイズによる失敗は各30件台にとどまり，無加工の元画像等での失敗はわずか12件であった．
通常の向きで撮影された画像にはほぼ失敗しないことを意味する一方，向きが大きく崩れた入力への対策（学習側拡張の強化または向き補正前処理）が今後の課題である．

\begin{table}[htbp]
\centering
\caption{誤識別の拡張種別内訳（単一識別・全6条件の合計）}
\label{tab:err_by_aug}
\small
\begin{tabular}{lr}
\hline
拡張種別 & 誤識別件数 \\
\hline
回転 $-$30度 & 109 \\
水平反転 & 108 \\
回転 $+$30度 & 89 \\
回転 $-$15度 & 63 \\
回転 $+$15度 & 50 \\
コントラスト 0.7倍 & 39 \\
コントラスト 1.3倍 & 38 \\
明度 0.7倍 & 35 \\
明度 1.3倍 & 34 \\
ノイズ & 30 \\
元画像・その他 & 12 \\
\hline
\end{tabular}
\end{table}

\textbf{（2）誤答先の偏り（受け皿薬剤）}：
よく起きる取り違えの組（正解→予測）の上位を表\ref{tab:confusion_pairs}に示す．

\begin{table}[htbp]
\centering
\caption{頻出する取り違えの組（単一識別・全6条件の合計・上位10組）}
\label{tab:confusion_pairs}
\small
\begin{tabular}{lr}
\hline
取り違え（正解 → 予測） & 件数 \\
\hline
パルモディアXR錠0.2mg → デエビゴ錠5mg & 14 \\
酸化マグネシウム錠330mg → プロマックD錠75mg & 13 \\
エンレスト錠100mg → プロマックD錠75mg & 12 \\
ミヤＢＭ錠 → バクトラミン配合錠 & 11 \\
ロスバスタチン錠2.5mg「DSEP」 → メインテート錠2.5mg & 11 \\
ジャヌビア錠50mg → デエビゴ錠5mg & 10 \\
ダイアート錠30mg → ジャヌビア錠50mg & 10 \\
メチコバール錠500$\mu$g → メトグルコ錠250mg & 10 \\
ノルバスクOD錠5mg → プロマックD錠75mg & 9 \\
カロナール錠500mg → ロキソプロフェンナトリウム錠60mg & 9 \\
\hline
\end{tabular}
\end{table}

誤答は特定の薬剤に偏って流れる傾向が見られた．
特にデエビゴ錠5mgとプロマックD錠75mgが複数の薬剤からの誤答先になっており，複数識別（検出→単剤分類）の誤検出でもデエビゴ錠5mgが81件と突出していた（2位のクエン酸第一鉄Na錠50mgは20件）．
これらの「受け皿」となる薬剤は白色円形の錠剤であり，外観の手がかりが乏しい場合の既定の答えになっていると考えられる．
また複数識別の取りこぼし（写っているのに当てられなかった薬剤）は，カロナール錠500mg（29件），ロスバスタチン錠2.5mg「DSEP」（22件），エンレスト錠100mg・ロキソプロフェンナトリウム錠60mg（各21件）が上位であり，単一識別の苦手薬剤と重なる．

\textbf{（3）同名・別用量の取り違え}：
エンレスト錠100mg→エンレスト錠200mg（3件），タケキャブ錠10mgと20mgの相互取り違え（各1件）など，同一製品名で用量のみが異なる薬剤間の取り違えが観察された．
件数は少ないが投与量の誤りに直結するため，医療応用上は最も注意すべき誤りであり，用量表示の読み取りを強化する対策が必要である．

\textbf{（4）量子化・小型化による誤りの増加}：
表\ref{tab:err_by_model}のとおり，単一識別の誤りは4B非量子化の27件に対し，0.8B量子化では209件と桁が変わる水準で増加した．
端末搭載時には，上記の苦手薬剤・受け皿薬剤を重点的に確認する必要がある．

\subsection{考察}

本章の実験結果から，以下の知見が得られた．

第一に，薬剤画像に特化したLoRAファインチューニングの効果は決定的である．
同一の評価画像814枚に対して，ファインチューニング前0.9\%に対しファインチューニング後98.2\%（差97.3ポイント）であり，汎用大規模モデルであるGemini（82.3\%）も上回った．
1薬剤あたり元画像3〜10枚という少数データでも，11種のデータ拡張との組み合わせで実用的な精度に到達できることが示された．
データ拡張なしの初期学習が過学習の挙動を示したこと（\ref{sec:noaug}節）を踏まえると，少数画像の薬剤鑑別においてデータ拡張は本手法の成立に不可欠な要素である．

第二に，複数薬剤の同時識別は，モデルの再学習ではなく「解き方の分解」で実現できた．
一括質問方式はプロンプトをどのように工夫しても0\%であったが，検出→単剤分類の2段構成により再学習なしで68.8\%に達し，さらに検出器の専用学習と入力解像度の向上（1024px）により96.0\%まで改善した．
検出器学習後の残差はすべて識別段階の取り違えであり，検出の問題は実質的に解消された．
一方で過検出抑制の後処理（NMS・該当なしゲート）は効果がなく，検出器そのものの学習が本質的な対策であったことが負の結果からも裏付けられた．

第三に，軽量化のトレードオフが定量化された．
単一識別のみであれば小型モデル（量子化2Bで83.9\%）も実用域に近いが，複数識別を含めるなら4Bが必須であり，量子化4B（単一91.9\%・複数55.2\%）が端末搭載の現実的な選択肢である．

最後に本実験の限界を述べる．
評価画像は（外部画像での予備的検証を除き）学習画像と同一環境で撮影したものであり，撮影環境の異なる画像（Web収集画像等）での汎化性能は未検証である．
複数薬剤の評価は合成画像によるものであり，実際に複数のPTPシートを並べて撮影した実写画像での検証も今後必要である．
また，回転・反転への脆弱性，同名・別用量の取り違えは未解決の課題として残る．

\newpage

%----------------------------------------------------------------------
\section{結論・今後の課題}\label{sec5}
%----------------------------------------------------------------------

\subsection{結論}

本研究では，オープンソースのVision-Language ModelであるQwen3.5-4Bを薬剤画像でLoRAファインチューニングし，薬剤画像とテキストプロンプトから薬剤名を直接出力する薬剤鑑別手法を提案・評価した．
得られた主な結論は以下のとおりである．

\begin{enumerate}
  \item \textbf{特化学習の有効性}：41薬剤・評価画像814枚の単一薬剤識別において，提案手法はTop-1精度98.2\%を達成し，ファインチューニング前のモデル（0.9\%）および外部API型マルチモーダルAIのGemini（82.3\%）を上回った．1薬剤あたり元画像3〜10枚という少数データでも，11種のデータ拡張との組み合わせにより高精度が実現できる．
  \item \textbf{複数薬剤同時識別の実現}：単剤識別用に学習したモデルを変更せず，「検出→単剤分類」の2段構成により複数薬剤の同時識別を実現した（完全一致率68.8\%，Gemini 50.6\%を上回る）．さらに検出器LoRAの専用学習と入力解像度の向上により96.0\%まで改善し，検出段階の課題を実質的に解消した．
  \item \textbf{軽量化のトレードオフ}：モデルサイズ（0.8B/2B/4B）と4bit量子化の影響を単一・複数・混在の3条件で定量化した．複数識別にはモデルサイズの影響が極めて大きく，端末搭載には量子化4B（単一91.9\%・複数55.2\%）が精度と軽さのバランスで最有力である．
  \item \textbf{ローカル動作}：以上の全評価はGemini比較を除きローカル環境（単一GPU）で完結しており，薬剤画像を外部サービスへ送信しない薬剤鑑別システムの実現可能性を示した．
\end{enumerate}

\subsection{今後の課題}

\begin{itemize}
  \item \textbf{別環境画像での汎化評価}：本研究の評価画像は学習画像と同一環境での撮影が中心である．Webから収集した別環境・別個体の画像で全薬剤の評価を行い，撮影環境に依存しない汎化性能を検証する（別時期撮影のレバミピド錠100mgでは99.7\%を確認済み）．
  \item \textbf{実写の複数薬剤画像での検証}：複数薬剤の評価は合成画像によるものであるため，実際に複数のPTPシートを並べて撮影した実写画像で最終検証を行う．
  \item \textbf{向きの崩れへの対策}：誤識別の66\%が回転・水平反転に起因した．学習側での大角度回転拡張の強化，または推論前の向き補正前処理を導入し，効果を比較する．
  \item \textbf{同名・別用量の取り違え対策}：エンレスト錠100mgと200mgのような用量違いの取り違えは投与量の誤りに直結する．用量表示の読み取りを強化する学習・プロンプト設計を検討する．
  \item \textbf{軽量モデルの拡充と実機検証}：Qwen3.5-2Bの再学習による「再学習あり2B」対「量子化のみ4B」の比較，およびスマートフォン実機でのオンデバイス推論の検証を行う．
  \item \textbf{RAG型手法との比較}：薬剤データベース検索を併用するRAG型手法（予備実験では5/5を確認）を全薬剤に拡張し，新規薬剤追加時に再学習が不要であるという利点を含めてファインチューニング型と比較する．
\end{itemize}

\newpage
\section*{謝辞}
\addcontentsline{toc}{section}{謝辞}

\newpage

\addcontentsline{toc}{section}{\bibname}
\bibliography{reference}
\bibliographystyle{junsrt}

\end{document}
